% Part: counterfactuals
% Chapter: minimal-change-semantics
% Section: antecedent-strengthening

\documentclass[../../../include/open-logic-section]{subfiles}

\begin{document}

\olfileid{cnt}{min}{agg}

\olsection{પૂર્વપક્ષને મજબૂત કરવો}

``પૂર્વપક્ષને મજબૂત કરવો'' એટલે
$!A \lif !C \Entails (!A \land !B) \lif !C$ એ અનુમાન.
ભૌતિક શરતી વિધાન માટે તે માન્ય છે, પરંતુ પ્રતિતથ્યાત્મક વિધાન માટે
અમાન્ય છે. ધારો કે આ દીવાસળી ઘસું તો તે સળગી ઊઠે, એ સાચું છે.
(એટલે દીવાસળી કે દીવાસળીઓના પૅકેટની ઘસવાની સપાટીમાં કોઈ ખામી નથી;
હું દીવાસળી તોડી નાખવાનો નથી, વગેરે.) પરંતુ આ દીવાસળી અવકાશમાં
ઘસું તો તે સળગી ઊઠે, એ સાચું નથી. તેથી નીચેનું અનુમાન અમાન્ય છે:
\begin{quote}
  જો દીવાસળી ઘસાત, તો તે સળગી ઊઠત.

  તેથી, જો દીવાસળી અવકાશમાં ઘસાત, તો તે સળગી ઊઠત.
\end{quote}

લૂઈસ--સ્ટાલનેકરનું શરતી વિધાનોનું વિશ્લેષણ આ સમજાવે છે: હું
અવકાશમાં દીવાસળી ઘસું એવું સૌથી નજીકનું જગત, હું દીવાસળી ઘસું
એવા સૌથી નજીકના જગત કરતાં વાસ્તવિક જગતથી ઘણું વધુ દૂર છે. તેથી
પછીના જગતમાં દીવાસળી સળગે છે, પરંતુ પહેલાંના જગતમાં નહિ. એવું જ
થવું જોઈએ.

\begin{ex}
  ગોળા-અર્થવિચાર આ અનુમાનને અમાન્ય ઠરાવે છે; એટલે $p
  \cif r \Entails/ (p \land q) \cif r$. નિદર્શ $\mModel{M}
  = \tuple{W, O, V}$ લો જેમાં $W = \{w, w_1, w_2\}$,
  $O_w = \{\{w\}, \{w, w_1\}, \{w, w_1, w_2\}\}$,
  $V(p) = \{w_1, w_2\}$, $V(q) = \{w_2\}$ અને
  $V(r) = \{w_1\}$. $p$ને સ્વીકારતો ગોળો $S = \{w,
  w_1\}$ છે અને તેના બધાં જગતોમાં $p \lif r$ સત્ય છે; તેથી
  $\mSat{M}{p \cif r}[w]$. વળી $(p \land q)$ને સ્વીકારતો ગોળો
  $S' = \{w, w_1, w_2\}$ છે, પરંતુ
  $\mSat/{M}{(p \land q) \lif r}[w_2]$; તેથી
  $\mSat/{M}{(p \land q) \cif r}[w]$ (જુઓ~\olref{fig:antecedent-strengthening}).

\begin{figure}
\begin{center}
\begin{tikzpicture}[layerwidth=1.5,scale=.6]\tiny
  \spheresystem{3}
  \propositionintersect{-10}{3}{30}{5}
  \begin{scope}
    \clip plot[smooth,tension=1] coordinates
          {(0,4.5) (30:.95) (4.5,-3.5)} -- (4.5,4.5) ;
    \spherefill{2}
  \end{scope}
  \draw[proposition] plot[dashed,smooth,tension=1] coordinates
          {(0,4.5) (30:.95) (4.5,-3.5)} ;
  \proposition{30}{2}{30}{5}
  \path (0,0) node[world] {$w$};
  \spherepos{30}{2}{node[world] {$w_1$}}
  \spherepos{-10}{3}{node[world] {$w_2$}}
  \spherepos{-10}{4}{node {$q$}}
  \spherepos{30}{4}{node {$r$}}
  \draw (1,4.5) node[below] {$p$};
\end{tikzpicture}
\caption{પૂર્વપક્ષને મજબૂત કરવાના નિયમનું પ્રત્યુદાહરણ}
\ollabel{fig:antecedent-strengthening}
\end{center}
\end{figure}
\end{ex}

\end{document}
